\section{BitIO Module}

\subsection{Purpose}

\texttt{bitio.hpp} provides a pair of thin wrappers around \texttt{std::ostream} and
\texttt{std::istream} that allow individual bits or groups of bits to be written and read.
It is the lowest-level primitive in the codebase: all LZSS token emission and decoding
goes through \texttt{BitWriter} and \texttt{BitReader}.

Without this abstraction, LZSS would need to manage its own bit packing --- mixing algorithm
logic with serialization detail.
By delegating that concern to a dedicated class, \texttt{lzss.cpp} can express its logic
in clean terms (\texttt{write\_bits(offset, 12)}) with no manual shifting.

\subsection{Files}

\texttt{include/bitio.hpp} only.
Both classes are fully defined in the header (all methods are short and inlineable),
so no \texttt{.cpp} file is needed.
This keeps the build simple: any translation unit that includes \texttt{bitio.hpp}
gets the full implementation at zero link cost.

\subsection{BitWriter}

\texttt{BitWriter} buffers bits until a full byte is available, then emits it to the
underlying stream.

\begin{lstlisting}[style=cpp]
class BitWriter {
public:
    explicit BitWriter(std::ostream& out);
    void write_bit(bool b);
    void write_bits(uint32_t val, int n);  // n LSBs of val, MSB first
    void flush();
private:
    std::ostream& out_;
    uint8_t buf_;         // accumulation byte (bits fill from MSB down)
    int bits_in_buf_;     // 0..7: bits currently in buf_
};
\end{lstlisting}

\paragraph{State}
\texttt{buf\_} accumulates bits left-to-right (MSB first).
\texttt{bits\_in\_buf\_} counts how many bits have been placed in \texttt{buf\_} since
the last byte was emitted.

\paragraph{write\_bit(b)}
Shifts \texttt{buf\_} left by one and ORs the new bit into the LSB position, then
increments \texttt{bits\_in\_buf\_}.
When the counter reaches 8 the byte is emitted and both fields are reset.

\paragraph{write\_bits(val, n)}
Iterates from bit \texttt{n-1} down to 0, calling \texttt{write\_bit} for each.
This writes the \texttt{n} least-significant bits of \texttt{val} with the most-significant
of those bits first (big-endian bit order within the group).

\paragraph{flush()}
If any bits remain in \texttt{buf\_} (\texttt{bits\_in\_buf\_} $>$ 0), the partial byte is
left-shifted to fill the remaining positions with zero bits, then emitted.
\textbf{The caller must call \texttt{flush()} after the last write.}
Forgetting this call silently drops the final partial byte.

\paragraph{Example: building byte \texttt{0xB0}}
Four calls \texttt{write\_bit(1)}, \texttt{write\_bit(0)}, \texttt{write\_bit(1)},
\texttt{write\_bit(1)} build \texttt{buf\_} to $\mathtt{0b00001011}$ (4 bits, pending).
A subsequent \texttt{flush()} left-shifts by 4 and emits \texttt{0xB0 = 0b10110000}.

\begin{center}
\begin{tabular}{r l r}
\toprule
\textbf{After call} & \textbf{\texttt{buf\_} (binary)} & \textbf{\texttt{bits\_in\_buf\_}} \\
\midrule
write\_bit(1) & \texttt{00000001} & 1 \\
write\_bit(0) & \texttt{00000010} & 2 \\
write\_bit(1) & \texttt{00000101} & 3 \\
write\_bit(1) & \texttt{00001011} & 4 \\
flush()       & emits \texttt{0xB0} (\texttt{10110000}) & 0 \\
\bottomrule
\end{tabular}
\end{center}

\subsection{BitReader}

\texttt{BitReader} fetches bytes from the stream on demand and returns bits one at a time.

\begin{lstlisting}[style=cpp]
class BitReader {
public:
    explicit BitReader(std::istream& in);
    bool read_bit();
    uint32_t read_bits(int n);  // reads n bits MSB-first
    bool eof() const;
private:
    std::istream& in_;
    uint8_t buf_;      // current byte being consumed
    int bits_left_;    // how many bits remain in buf_
};
\end{lstlisting}

\paragraph{read\_bit()}
When \texttt{bits\_left\_} reaches 0, the next byte is fetched from the stream into
\texttt{buf\_} and \texttt{bits\_left\_} is set to 8.
The most-significant remaining bit is extracted as
\texttt{(buf\_ >> bits\_left\_) \& 1}, and \texttt{bits\_left\_} is decremented.
On stream EOF the method returns \texttt{false}.

\paragraph{read\_bits(n)}
Accumulates \texttt{n} calls to \texttt{read\_bit()} into a \texttt{uint32\_t},
shifting the accumulator left after each bit.
This reconstructs the MSB-first value that \texttt{write\_bits} encoded.

\paragraph{eof()}
Returns \texttt{true} when the underlying stream is exhausted and no bits remain
in \texttt{buf\_}.

\subsection{Design Notes}

\textbf{No ownership.}
Both classes hold a \emph{reference} to the stream, not a pointer or \texttt{unique\_ptr}.
The caller owns and manages the stream's lifetime; BitWriter/Reader are lightweight
view objects that must not outlive the stream they wrap.

\textbf{Header-only.}
All methods are short enough that the inline cost is negligible.
Keeping both classes in a single header avoids an extra translation unit and
simplifies the build.
